% AY2020 制御工学 解答例
% 回答の流れ・言葉遣いは 03_制御工学/参考/「制御工学Ⅱ（完成版）」に寄せる。
% デザイン（囲み枠等）は真似しない。解答・数値は本年度過去問から自力で導いた。
\documentclass[11pt,a4paper]{article}
\usepackage{../../../00_common/preamble}

\usetikzlibrary{positioning,arrows.meta,calc}
\usepackage{pgfplots}
\pgfplotsset{compat=newest}

\title{制御工学 AY2020 解答例（専門応用科目I）}
\author{}
\date{}

\begin{document}
\maketitle

% ==================================================================
\section*{第1問〔I〕}

台車（質量 $m$、変位 $x_2$）と入力端 $x_1$ の間に、直列バネ $k_1$-$k_2$ とダンパ $d$ を
並列に接続した系。直列バネの合成は
\[
  k_{\mathrm{eq}}=\frac{k_1k_2}{k_1+k_2}.
\]

\subsection*{(1) 運動方程式と伝達関数}
バネ（合成 $k_{\mathrm{eq}}$）とダンパ $d$ がともに質量と入力端 $x_1$ を結ぶので
\begin{align*}
  &\qquad m\ddot{x}_2=k_{\mathrm{eq}}(x_1-x_2)+d(\dot{x}_1-\dot{x}_2)\\
  &\therefore\quad m\ddot{x}_2+d\dot{x}_2+k_{\mathrm{eq}}x_2=d\dot{x}_1+k_{\mathrm{eq}}x_1\\
  &\therefore\quad (ms^2+ds+k_{\mathrm{eq}})X_2=(ds+k_{\mathrm{eq}})X_1.
\end{align*}
\[
  \therefore\quad \ans{\;G(s)=\dfrac{ds+k_{\mathrm{eq}}}{ms^2+ds+k_{\mathrm{eq}}},\qquad k_{\mathrm{eq}}=\dfrac{k_1k_2}{k_1+k_2}\;}
\]

\subsection*{(2) ステップ応答と最大値（$m=1,d=2,k_1=k_2=2$）}
$k_{\mathrm{eq}}=\dfrac{2\cdot2}{4}=1$ より $G(s)=\dfrac{2s+1}{s^2+2s+1}=\dfrac{2s+1}{(s+1)^2}$。$X_1=\dfrac1s$ なので
\begin{align*}
  &\qquad X_2(s)=\frac{2s+1}{s(s+1)^2}=\frac1s-\frac{1}{s+1}+\frac{1}{(s+1)^2}\\
  &\therefore\quad x_2(t)=1-\e^{-t}+t\e^{-t}=1-(1-t)\e^{-t}.
\end{align*}
最大値は $\dot{x}_2=(2-t)\e^{-t}=0$ より $t=2$、そのとき $x_2=1+\e^{-2}$。
\[
  \therefore\quad \ans{\;x_2(t)=1-(1-t)\e^{-t},\qquad x_{2\max}=1+\e^{-2}\approx1.14\ (t=2)\;}
\]

\subsection*{(3) 合成入力応答（$m=1,d=2,k_1=2,k_2=4$）}
$k_{\mathrm{eq}}=\dfrac{2\cdot4}{6}=\dfrac43$ より $G(s)=\dfrac{2s+\frac43}{s^2+2s+\frac43}$。極は $s=-1\pm\dfrac{j}{\sqrt3}$。
$X_1(s)=\dfrac1s+\dfrac{1}{s+1}$（$x_1=1(t)+\e^{-t}$）として $X_2=G X_1$ を逆変換すると
\[
  \therefore\quad \ans{\;x_2(t)=1-2\e^{-t}+\e^{-t}\!\left(\cos\tfrac{t}{\sqrt3}+3\sqrt3\,\sin\tfrac{t}{\sqrt3}\right)\;}
\]
減衰振動項は $\e^{-t}$ で消える。充分な時間が経過すると入力は単位ステップ成分だけが残り、
$G(0)=1$ より
\[
  \therefore\quad \ans{\;\lim_{t\to\infty}x_2(t)=1\;}
\]

\subsection*{(4) 無減衰の場合（$m=2,d=0,k_1=1,k_2=2$）}
$k_{\mathrm{eq}}=\dfrac{1\cdot2}{3}=\dfrac23$、$d=0$ より
\begin{align*}
  &\qquad G(s)=\frac{2/3}{2s^2+2/3}=\frac{1/3}{s^2+1/3}\\
  &\therefore\quad X_2(s)=\frac{1/3}{s\,(s^2+\frac13)}=\frac1s-\frac{s}{s^2+\frac13}\\
  &\therefore\quad x_2(t)=1-\cos\frac{t}{\sqrt3}.
\end{align*}
\[
  \therefore\quad \ans{\;x_2(t)=1-\cos\dfrac{t}{\sqrt3}\;}
\]
$d=0$ なので減衰せず、$1$ を中心に振幅 $1$ で持続振動する。

\bigskip
\noindent 検算: (2) は $x_2(0)=0$・$t=2$ で $x_{2\max}=1+\e^{-2}$、(3) は減衰振動項が消えて定常 $1$、
(4) は極 $\pm j/\sqrt3$ の無減衰振動を SymPy で確認✓。

\newpage
% ==================================================================
\section*{第2問〔II〕}

前向き $G_1G_2$、帰還 $K$ のフィードバック系。$G_1=\dfrac{K_1}{s+1}$、$G_2=\dfrac{1}{0.1s+1}$。

\subsection*{(1) 一巡伝達関数と単位ステップ応答（$K_1=2,K=1$）}
一巡伝達関数はループ一周の積 $L(s)=K G_1 G_2$。
\begin{align*}
  &\qquad L(s)=1\cdot\frac{2}{s+1}\cdot\frac{1}{0.1s+1}=\frac{2}{(s+1)(0.1s+1)}.
\end{align*}
$s=j\omega$ として
\[
  \therefore\quad \ans{\;|L(j\omega)|=\dfrac{2}{\sqrt{1+\omega^2}\,\sqrt{1+(0.1\omega)^2}},\quad
  \angle L(j\omega)=-\tan^{-1}\omega-\tan^{-1}0.1\omega\;}
\]
閉ループは
\begin{align*}
  &\qquad T(s)=\frac{G_1G_2}{1+KG_1G_2}
   =\frac{2}{(s+1)(0.1s+1)+2}=\frac{20}{(s+5)(s+6)}.
\end{align*}
$X(s)=\dfrac{T(s)}{s}$ を分解して
\[
  \therefore\quad \ans{\;x(t)=\dfrac23-4\e^{-5t}+\dfrac{10}{3}\e^{-6t}\;}\qquad(\text{定常値}\ \tfrac23)
\]

\subsection*{(2) ゲイン線図（漸近線）}
$|L(0)|=2$、すなわち $20\log_{10}2\approx6\,\mathrm{dB}$。折れ点は $s+1$ より $\omega=1$、$0.1s+1$ より
$\omega=10\,[\mathrm{rad/s}]$。傾きは $\omega<1$ で $0$、$1<\omega<10$ で $-20$、$\omega>10$ で $-40\,[\mathrm{dB/dec}]$。
（(4) の場合を破線で併記。）

\begin{center}
\begin{tikzpicture}
\begin{axis}[
  width=0.86\linewidth, height=5.0cm,
  xmin=1e-2, xmax=1e3, ymin=-60, ymax=20, xmode=log,
  xlabel={$\omega\,[\mathrm{rad/s}]$}, ylabel={ゲイン [dB]},
  ytick={-60,-40,-20,0,6,20},
  xmajorgrids=true, ymajorgrids=true, xminorgrids=true,
  grid style={line width=0.3pt, draw=gray!60},
  extra y ticks={0}, extra y tick style={grid=major, grid style={line width=1pt, draw=black}},
  legend pos=south west, clip=true, enlargelimits=false,
]
% (2): K=1, DC 6dB, 折れ点 1,10
\addplot[thick] coordinates {(0.01,6) (1,6) (10,-14) (1000,-94)};
\addlegendentry{(2) $K{=}1,\,G_2{=}\frac{1}{0.1s+1}$}
% (4): K=0.5, G2=1/(0.2s+1), DC 0dB, 折れ点 1,5
\addplot[thick,dashed] coordinates {(0.01,0) (1,0) (5,-14) (1000,-60)};
\addlegendentry{(4) $K{=}0.5,\,G_2{=}\frac{1}{0.2s+1}$}
\node[anchor=south] at (axis cs:1,6) {\scriptsize $1$};
\node[anchor=south west] at (axis cs:10,-14) {\scriptsize $10$};
\node[anchor=north] at (axis cs:5,-14) {\scriptsize $5$};
\end{axis}
\end{tikzpicture}
\end{center}

\subsection*{(3) ゲインを $0$ dB にする $K$}
$\omega=0.01$ は折れ点 $\omega=1$ より充分低く、$|L|$ は低周波の直流ゲイン $|L(0)|=K\,K_1=2K$ に等しい。
これを $0\,\mathrm{dB}=1$ とおくと
\[
  2K=1\quad\therefore\quad \ans{\;K=\dfrac12\;}
\]

\subsection*{(4) $G_2(s)=\dfrac{1}{0.2s+1}$、$K=\tfrac12$ の場合の比較}
一巡伝達関数は
\begin{align*}
  &\qquad L(s)=\frac12\cdot\frac{2}{s+1}\cdot\frac{1}{0.2s+1}=\frac{1}{(s+1)(0.2s+1)}.
\end{align*}
$|L(0)|=1=0\,\mathrm{dB}$、折れ点 $\omega=1,\ 5$。閉ループは
\begin{align*}
  &\qquad T(s)=\frac{G_1G_2}{1+KG_1G_2}
   =\frac{2}{(s+1)(0.2s+1)+1}=\frac{10}{s^2+6s+10}=\frac{10}{(s+3)^2+1}.
\end{align*}
極は $s=-3\pm j$（不足減衰、$\zeta=\dfrac{3}{\sqrt{10}}\approx0.95$）。ステップ応答は
\[
  \therefore\quad \ans{\;x(t)=1-\e^{-3t}\bigl(\cos t+3\sin t\bigr)\;}\qquad(\text{定常値}\ 1)
\]
\textbf{考察}：
\begin{itemize}\setlength{\itemsep}{1pt}
  \item \textbf{ゲイン線図}：$K$ を $1\to\tfrac12$ にしたので低周波ゲインが $6\to0\,$dB に下がり、
        $G_2$ の折れ点が $\omega=10\to5$ に移って $-20\,$dB/dec の帯域が狭まる。
  \item \textbf{時間応答}：(2) の系（$G_2=\frac{1}{0.1s+1}$）の閉ループ極は $-5,-6$ の実根（過減衰、
        オーバーシュートなし）だが、(4) では $G_2$ の時定数が大きくなり閉ループ極が $-3\pm j$ の
        複素根（不足減衰）となるため、わずかにオーバーシュートを伴う応答に変わる。
  \item 定常値は (1) の $\tfrac23$ から $K$ 調整により (4) では $1$（$0\,$dB に対応）に変化する。
\end{itemize}

\bigskip
\noindent 検算: (1) $T=20/((s+5)(s+6))$ 定常 $\tfrac23$、(3) $K=\tfrac12$、
(4) $T=10/((s+3)^2+1)$ 極 $-3\pm j$・定常 $1$ を SymPy で確認✓。

\end{document}
